\documentclass{article}
\usepackage[T2A]{fontenc}
\usepackage[utf8]{inputenc}
\usepackage{amsthm}
\usepackage{amsmath}
\usepackage{amssymb}
\usepackage{amsfonts}
\usepackage{mathrsfs}
\usepackage[12pt]{extsizes}
\usepackage{fancyvrb}
\usepackage{indentfirst}
\usepackage[
  left=2cm, right=2cm, top=2cm, bottom=2cm, headsep=0.2cm, footskip=0.6cm, bindingoffset=0cm
]{geometry}
\usepackage[english,russian]{babel}


\begin{document}
\section*{Вариант 21}

Зная зависимости изменения плотности газа при нагнетании, мы можем рассчитать производительность компрессора как расход газа через третью ступень компрессора:

\begin{equation}
  G(t) = \lambda_3 V_{\Pi3} n \rho_3 (t)
\end{equation}
\noindent где $\rho_3 (t)$ — известная зависимость плотности газа от времени на 3-й ступени компрессора.

Теперь мы можем рассчитать потребление электроэнергии компрессором:

\begin{align}
  N(t) &= G(t) R \frac{k}{k - 1} \sum^{2}_{j = 1} (\vartheta_{Hj}(t) - \vartheta_0) + \Delta N, \\
  \Delta N &=
  \frac{G(t) * R * T_1 * \delta P_1}
       {\frac{V_{\Pi1}}{V_{\Pi2}} P_{ATM}} +
  \frac{G(t) R T_2 \delta P_2}
       {(\frac{V_{\Pi1}}{V_{\Pi2}} P_{ATM} - \delta P_1) \frac{V_{\Pi2}}{V_{\Pi3}}}
\end{align}

\noindent где $G(t)$ "--- искомая производительность компрессора; $\Delta N$ "--- потери мощности в промежуточных газоохладителях; $R$ "--- газовая постоянная; $T_1$, $T_2$ "--- среднее арифметическое температур воды и воздуха в воздухоохладителях; $\delta P_1$, $\delta P_2$ "--- потери давления воздуха в воздухоохладителях; $P_{ATM}$ — атмосферное давление.

\end{document}
